\begin{abstract}
Attribution methods are widely used to explain which sparse-autoencoder (SAE) features drive a probe's prediction, but it is rarely tested whether those attributions predict what happens when the features are changed in the language model. We attribute linear and MLP sentiment probes on GPT-2 layer-7 SAE features with SHAP, Integrated Gradients and Gradient$\times$Input, and compare each to the measured logit-difference effect of steering or ablating the same features on a sentence-level SST-2 split. \textbf{On direction, the probe's raw weights beat signed global SHAP: they predict the sign of a feature's effect ($\rho=0.31$ last-token, $0.40$ all-token), while SHAP, which starts from those weights and adjusts them for the data, does not ($\rho=-0.11$ and $-0.10$; gap $0.41$, 95\% CI $[0.09,0.72]$).} The failure comes from centring the signed attribution on a background mean, not from the probe: SHAP is anti-correlated with the probe's own closed-form response, and a background-free score recovers the direction. For ranking feature importance, SHAP, IG and GA, aggregated the same way, are modestly faithful ($\rho=0.42$--$0.50$) and nearly indistinguishable from one another, and in the linear arm no better than a background-free product of weight and activation. Local SHAP outranks global SHAP only when inactive features are included, where it benefits from knowing which are active; among active features we find no advantage. The design cannot separate attribution error from probe--language-model mismatch beyond the moderate link between them ($\rho\approx0.3$--$0.4$), so we claim no unfaithfulness to the language model beyond that level.
\end{abstract}
